% page 250 of the facsimile (image p-249 of the scan)
% block 160.0 x 254.9 mm ; left margin 8.0 mm
\pgc{-12.700mm}{0.000mm}{
\l{242}
\l{}
\l{\sou{invektivar}. (\sou{netrans., kontre}). Diskursar violente, furio-}
\l{\cel{3}ze,quankam mem decoze, kontre ulu od ulo. - DEFIS.}
\l{}
\l{\sou{inventariar}. (\sou{trans.)}\cel{1}Kontar ed enrejistrar specope la hava-}
\l{\cel{3}ji transportebla,vari, komerco-trati, mandati, acioni,}
\l{\cel{3}obligacioni di firmo, di persono, e c.) - DEFIRS.}
\l{}
\l{\sou{inventar}. (\sou{trans.})\cel{2}I. Krear ulo nova. - II. (\sou{retoriko)}.}
\l{\cel{3}Serchar e selektar la argumento-punti quin onu devas uzar,}
\l{\cel{3}la idei quin la temo furnisas. - EFIRS.}
\l{}
\l{\sou{inversa}.\cel{2}Qua esas, o qua venas, opozo-sinse. - EFIS.}
\l{}
\l{\sou{invertazo}. (\sou{kemio)}.\cel{2}Fermento solvebla, qua transformas sako-}
\l{\cel{3}roso aden glikoso,e qua existas en vejetanti diversa.}
\l{}
\l{\sou{investar}.\cel{2}(\sou{trans.})\cel{2}I. Solene igar posedeskar feudo, granda}
\l{\cel{3}ofico, per ceremonio alegoriala. - II. Igar posedeskar}
\l{\cel{3}povo, autoritato. - DEFIS.}
\l{}
\l{\sou{invetera}.\cel{2}Qua lasis kreskar en su - til divenir neextirpe-}
\l{\cel{3}bla pro lua ancieneso - kustumo mala. (ex.: lu esas fumero,}
\l{\cel{3}hiper-drinkero, invetera). EF.}
\l{}
\l{\sou{invitar}.\cel{2}(\sou{trans.,ulu, ad ulo, pri}). Pregar (ulu) ke lu irez}
\l{\cel{3}a loko determinita. - EFIS.}
\l{}
\l{\sou{invokar}. (\sou{trans.)}\cel{2}Advokar per prego ; advokar (ulu) por ke}
\l{\cel{3}lu atestez favore di la advokanto. EF.}
\l{}
\l{\sou{involuciono}. I. (\sou{geom.)}\cel{2}Dui de punti AA\textquotesingle{}+BB\textquotesingle{}+BB\textquotesingle{}+CC\textquotesingle{} - se-}
\l{\cel{3}lektita sur rekto, esas involuciono kande li dividas har-}
\l{\cel{3}monioze un segmento MN di la rekto. - II. (\sou{bot.}) Eso-ma-}
\l{\cel{3}niero di parto volvita adinterne di su. - DEFIRS.}
\l{}
\l{\sou{involukro}. (\sou{bot.)}\cel{2}Asemblajo ek braktei qui formacas, cirkum}
\l{\cel{3}floro o kapitulo, quaze kalico. - DEFIS.}
\l{}
\l{\sou{inyamo}. (\sou{bot.}) Planto tropikala di qua la stipito esas klimera,}
\l{\cel{3}ek la familio \textquotedbl{}amilacei\textquotedbl{}, e di qua la radikaro esas nutriva.}
\l{\cel{3}- L.\cel{1}\sou{discorea}. - DEFIRS.}
\l{}
\l{\sou{iodo}.\cel{2}(\sou{kemio)}. Korpo simpla, ek flitri briloza, bluatre-griza,}
\l{\cel{3}metalea, saporanta feble quale kloro. -Simb. kem.:I.- DEFIRS}
\l{}
\l{\sou{iodoformo}. (\sou{kemio)}.\cel{2}Kompozajo(\textquotedbl{}CHI3) deskovrita en 1824 da}
\l{\cel{3}Serullas.}
\l{}
\l{\sou{iono.}\cel{2}(\sou{fiziko}). Partikulo qua portas elektro pozitiva (iono}
\l{\cel{3}pozitiva) o negativa (iono negativa). - DEF.}
\l{}
\l{\sou{ionika}.\cel{2}(\sou{arkitekt.)}\cel{2}(\sou{aludante la triesma ek la kin klasA}}
\l{\cel{3}\sou{arkitekturala antiqua}). Karakterizata da la voluti di la}
\l{\cel{3}kapiteli. - DEFIS.}
\l{}
\l{\sou{\textquotedbl{}iota\textquotedbl{}}. La litero nonesma, e la maxim mikra, ek la alfabeto}
\l{\cel{3}Greka, qua korespondas ad \textquotedbl{}y\textquotedbl{}.}
}
